% !TeX encoding = UTF-8
% !TeX spellcheck = en_GB
% !TeX root = ../thesis.tex

\chapter{Related Work}
\label{chapter:related_work}
\section{PUFs on FPGAs}
Numerous PUFs suitable for implementation on FPGAs have been proposed. However, very few publications (two to our knowledge) provide the source code of their designs, limiting reproducibility, external evaluation, and practical reuse or extension in future works. As highlighted by \cite{6546024} and \cite{6645555}, one of the main motivations for implementing PUFs on FPGAs is the protection of hardware intellectual property (HW-IPs) by binding it to specific FPGA devices. The hardware IP protection consists of protecting the generated bitstream, which is the final compiled output from the hardware description source code that is used to program the FPGA. This bitstream is typically stored on insecure flash memory on the same board as the FPGA. If the bitstream is illegally copied, this can lead to revenue loss due to it being intellectual property. Such IP is very valuable since it is the main product of companies that build application specific high performance processors (ASHPP) using FPGAs. Previous binding techniques rely on trusted third parties for public-key encryption \cite{4439246} and store keys on secure ROM or flash memory, which exposes the IP to side-channel attacks. Researchers are therefore investigating the usage of PUFs on FPGAs as a safer and cheaper alternative for storing secrets critical to IP binding mechanisms.


\cite{stephan_nolting_2021_5706593} implemented and briefly tested a 96-bit ring oscillator PUF on various FPGA boards, but found that only 63 bits were stable enough for a reliable response. He also observed that repeated successive queries increased the likelihood of bit flips, which he attributes to a rise in chip temperature. While the high number of successive queries (4096 to be exact) can be justified when using response processing methods such as voting or hysteresis - where a bit is considered "stable" if it is consistently set (to 1) or cleared (to 0) above a given threshold, this introduces trade-offs between reliability, query speed, and thermal constraints when using FPGA PUFs in practical devices. The need for such a high number of measurements could be mitigated by using error-correcting codes (ECCs), which would allow occasional unstable bits to be corrected without the need for repeated measurements.
\cite{9556423} designed and implemented a new delay-based PUF called CHOICE PUF. The proposed PUF supports 6144 challenges but can be adapted to support more. One particularity of the CHOICE PUF is that it is tunable, meaning that they could increase or decrease the probability of obtaining 1's in the PUF response using the properties of the challenges. They also found that 147 of the 6144 challenges provided a uniformity rate between 49\% and 51\%. Furthermore, on the 147 challenges with desirable uniformity, the error rate went up to 4\% and as low as 1.5\%, making it a very reliable PUF. The same 147 challenges also had inter-PUF uniqueness near 50\%.
In an effort to reduce the footprint of PUFs on FPGAs for allowing the FPGA to be used for other applications while still housing a PUF, \cite{5547067} propose time division multiplexed PUFs (TDM-PUFs), which is a method to divide one long PUF into several smaller PUFs that run in different times by exploiting the dynamic partial reconfiguration (DPR) feature available in most FPGAs, which allows reconfiguring the FPGA during operation to use the same hardware resources for different purposes in different time segments.

As summarized by \cite{cryptography7040055}, whilst FPGA-based PUF designs provide advantages for device authentication and key generation, they also introduce important trade-offs. Notably, the reliability of PUF responses can be affected by environmental factors such as temperature. Also, the wide variety of existing FPGA-based PUF designs introduces challenges in terms of standardization and scalability, as the PUF implementations may require significant FPGA resources or careful placement and routing constraints, which can limit scalability when large CRPs are needed.

\section{Reconfigurable PUF}

Reconfigurable Physical Unclonable Functions (R-PUF) are PUFs whose challenge-response behavior can be updated. \cite{10.1145/1502781.1502786} first proposed leveraging FPGA reconfiguration to modify a PUF's behavior. \cite{10.1007/978-3-642-23951-9_25} later formalized the concept with Logically Reconfigurable PUFs (LR-PUFs), which use an internal state to remap challenges and responses without modifying the underlying hardware. Invoking a reconfiguration function replaces the old state with a new one (e.g: through hashing), producing a new CRP set. For a more practical implementation of an LR-PUF, \cite{LR-PUF-Memory-Based} introduce a memory-based LR-PUF that combines an SRAM PUF with a non-volatile state in NVM (non-volatile memory), the response of the PUF is concatenated with the current state and hashed to hide the PUF's behavior and protect it from potential modeling attacks. The key principle is that the new behavior of the PUF after reconfiguration is unpredictable from the previous one. This hiding of the PUF's behavior protects it against cloning or modeling attacks. 

\cite{10.1145/3274666} classified R-PUFs into two categories. The first is circuit-based R-PUFs, where the physical circuit is reconfigured (e.g: via FPGA routing) to produce a new hardware instance of the PUF. For example, \cite{DPUF} propose a reconfigurable DRAM PUF based on refresh pausing, where the reconfiguration is done by changing the refresh-pause interval to alter the CRPs generated by the PUF. The second category is algorithm-based R-PUFs, where the hardware remains the same but the challenge-to-response mapping is altered through extra logic. As proposed by \cite{10.1007/978-3-642-23951-9_25}, a transformation is applied to an input challenge using the current state, and similarly another transformation is applied to the output using the state. Here, the configuration of the PUF is simply comprised of extra steps that are applied on its output and inputs.

Reconfigurable PUFs are leveraged in security protocols where they are used in IoT scenarios to provide fresh keys and counter replay or modeling attacks. For example, \cite{10549865} propose a mutual authentication protocol where an R-PUF is reconfigured after each session to limit the exposure of its CRPs. \cite{s22155577} propose an R-PUF-based authentication framework for FPGA-based IoT devices, where the previously mentioned dynamic partial reconfiguration feature is used to reconfigure the PUF by changing its placement at runtime.

Because reconfiguration scrambles the CRP mapping each time, R-PUFs offer strong security advantages. \cite{10.1007/978-3-642-23951-9_25} proves forward- and backward-unpredictability of their LR-PUF model, i.e. responses measured before the reconfiguration are invalid after it, and an attacker with access to a reconfigured PUF cannot estimate the PUF behavior before reconfiguration.


\section{Methods Against Session Hijacking}
The most commonly used method for preventing session hijacking in web communications is the HttpOnly cookie attribute introduced in RFC6265 \cite{rfc6265}. The HttpOnly attribute prevents JavaScript from directly accessing session cookies, thereby preventing attacks in which cross-site scripting (XSS) is used to steal cookies. Another used attribute is the SameSite cookie attribute \cite{ietf_rfc6265bis_draft} which restricts where browsers send cookies to mitigate previous cases where sensitive cookies were unintentionally transmitted to malicious sites upon visit. While these mechanisms cover some attack vectors leading to session hijacking, the session cookies may still be compromised through other means such as malware \cite{tyler2025browsercontrolsprotectcookies, cookieMiner_unit42_2019} or network interception \cite{evilginx2_kgretzky_2024} as described in \cite{mitre_t1539}. As a result, several research efforts have explored methods to tackle session hijacking at a deeper level.

With the exception of the aforementioned PUSH protocol, no approach to defend against session hijacking in web communications uses PUFs, rather the usage of PUFs is mostly limited to IoT authentication schemes where they are used as a source of random numbers \cite{10.1007/978-3-540-88625-9_3}, unique identifiers \cite{9557280} or cryptographic keys \cite{10538739, 10007641}. 

The methods for protecting web communication or any other kind of communication against session hijacking usually involve binding a session to some other variable through continuous authentication (CA), where the user is repeatedly authenticated throughout the session using contextual signals rather than relying solely on the initial login; for instance, the variable could be user behavior \cite{6879297, 8757539}, user location \cite{5692487}, device characteristics \cite{9530399}, or biometrics \cite{https://doi.org/10.1002/widm.1365, 8010662, 6424854}. \cite{8757539} propose an adaptive CA method that models the keystroke patterns and mouse usage behavior patterns of a user to continuously verify identity during a session. \cite{9530399} propose an IoT CA protocol that involves devices features such as the battery level and device location. \cite{6424854} evaluate the usage of biometrics in CA and adapts the session timeouts based on user activity and biometric data quality. 

Other approaches focus on limiting the exposure of session cookies. For example, \cite{10.1007/978-3-642-19125-1_7} tackle the problem preemptively by limiting the exposure of the session cookies to XSS attacks through hiding the session cookies from the browser through an external proxy. Similarly, \cite{10.1007/978-3-319-64701-2_7} developed a proxy to detect and drop any illegitimate transfer of cookies. Device Bound Session Credentials (DBSC) \cite{w3c-dbsc} propose binding sessions to device-protected cryptographic keys, typically secured using hardware such as a Trusted Platform Module (TPM), instead of only relying on bearer cookies. 

A separate class of session protection mechanisms relies on one-time session tokens. Such protocols prevent replay attacks by generating a fresh token for each request. \cite{He2020} propose and evaluate a one-time cookies (OTC) protocol based on a novel hash-based one-time password (OTP) protocol, which unlike traditional OTP approaches does not rely on time synchronization or on a fixed shared key for an entire session, rather multiple hash chains are established during login to generate one-time values that encrypt and authenticate each subsequent request.

Each of the aforementioned approaches presents its own advantages and limitations. For some methods, the deployment often introduces additional complexity and infrastructure requirements, making them more suited for environments with high security requirements. Others are easier to integrate, but offer weaker protection. Some use AI or heuristics for continuous verification, which introduces false acceptance rates (FAR) and false rejection rates (FRR), whereas others are grounded on solid mathematically proven cryptographic protocols. As a result, defending against session hijacking is a trade-off between security, usability, and ease of deployment. This compromise is important to keep in mind when considering methods against session hijacking.
